\documentclass{situationist-preprint}

\title{The Absurd and the Appropriate Response}
\msid{EXI-04-a}
\msfamily{EXI-04}
\msdiscipline{Meaning, Death and the Absurd}
\mstype{Principal manuscript}
\mskeywords{absurd; Nagel; Camus; irony; despair; fitting attitudes; detached standpoint}

\begin{document}
\maketitle

\begin{abstract}
Is life absurd, and if so, how should one respond? \citet{Nagel1971} located the absurd not in the smallness or brevity of human life but in a collision within the reflective subject: we take our lives with a seriousness that we cannot help, and we can always step back and see that seriousness as arbitrary and open to doubt. We argue that life is absurd in this sense, because both of the colliding features are ineliminable for any reflective agent, and that the appropriate response is irony rather than despair. The key move is to treat responses to the absurd as attitudes assessable for fit. Despair presents the detached standpoint as having shown that our ends are worthless; it has shown no such thing, since the step back yields the absence of a vindication, not a refutation. Camus's revolt presents the world as an adversary to be scorned, which misdescribes an indifferent world. Irony presents a discrepancy between the seriousness of our commitments and the support available for them from the most detached standpoint, and that discrepancy is exactly what the absurd consists in. Irony is therefore the fitting response, and it is compatible with full engagement in worthwhile projects.
\end{abstract}

\section{The problem}

Many arguments for the absurdity of human life appeal to scale. We are small in an immense universe, our lives are brief in an immense stretch of time, and nothing we do will matter in a million years. \citet{Nagel1971} observed that none of these considerations does the work it is meant to do. If a life of seventy years is pointless, a life of eternity would be no less so; if our smallness makes us insignificant, being the size of the universe would not help. The considerations express something, but they do not state it. What they express, on Nagel's diagnosis, is a structural feature of reflective life: the collision between the seriousness with which we pursue our ends and the perpetual possibility of stepping back from them, to a standpoint from which they look arbitrary and from which no justification of them is final. Justifications within a life come to an end in ends that are not themselves justified, and reflection can always notice this.

Nagel added two claims about the response. Since the detached standpoint is one we cannot occupy permanently, and since what it reveals gives no reason for anything, the absurd does not warrant despair; and the heroic defiance recommended by \citet{Camus1942}, who described Sisyphus as conscious of his fate and superior to it through scorn, takes the cosmic standpoint more seriously than it deserves. Nagel's recommendation was that we approach our absurd lives with irony. \citet[ch.~XI]{Nagel1986} returned to the question and treated the collision as one instance of the general tension between the subjective and the objective standpoints.

Two lines of later discussion bear on the recommendation. One asks whether the absurd is a defect at all: if it follows from the capacity for reflection, it is the shadow of a capacity we have every reason to value. The other asks whether objective value, if there is any, dissolves the absurd: on theories of meaning such as those of \citet{Wolf2010} and \citet{Metz2013}, our projects may be worthwhile independently of our attitudes toward them, and then the step back seems to reveal nothing troubling. We argue that life is absurd in Nagel's sense even if those theories are true, and that irony is the attitude that fits it. Camus serves as our foil throughout.

\section{Framework}

\begin{definition}[Seriousness]\label{d:serious}
An agent takes an end \emph{seriously} if she devotes effort and attention to it, treats considerations bearing on it as reasons, and is liable to distress and satisfaction depending on how it fares.
\end{definition}

\begin{definition}[Detached standpoint]\label{d:detach}
The \emph{detached standpoint} on an agent's ends is the standpoint from which she considers those ends as the ends of one creature among others, without presupposing their importance, and asks what justifies her concern with them.
\end{definition}

\begin{definition}[Absurd]\label{d:absurd}
A life is \emph{absurd} if its subject takes some ends seriously and, from the detached standpoint, no justification of that seriousness is available that is not itself open to the same question.
\end{definition}

Definition~\ref{d:absurd} is Nagel's notion. It does not say that the ends are worthless, only that the regress of justification has no terminus that the detached standpoint must accept. It also does not say that the subject is mistaken, unhappy or foolish. Absurdity, so understood, is a relation between two attitudes of the same subject.

\begin{principle}[Fit]\label{p:fit}
An attitude toward a situation is \emph{fitting} to the extent that the evaluative features under which it presents the situation are features the situation has.
\end{principle}

Principle~\ref{p:fit} is the standard fitting-attitude idea applied to responses: fear presents its object as dangerous and is fitting when the object is dangerous; pride presents an achievement as one's own and creditable. The question of the appropriate response to the absurd is then the question of which attitude presents the absurd situation as it is.

\begin{definition}[Three responses]\label{d:resp}
\emph{Despair} toward one's ends presents them as worthless or as beyond the reach of any good. \emph{Revolt}, in Camus's sense, presents the world as an adversary that denies us meaning, and the agent's persistence as a victory over it. \emph{Irony} toward one's ends presents a discrepancy between the seriousness with which one holds them and the support for that seriousness available from the most detached standpoint, while continuing to hold them.
\end{definition}

The irony in question is not the rhetorical figure of saying the opposite of what one means, nor the posture of noncommitment. It is the attitude of someone who sees that her commitment claims more than can be vindicated from outside and keeps the commitment without pretending otherwise. It resembles the attitude of a player who knows that the game's stakes exist only within the game and still plays to win.

\section{The argument}

\begin{lemma}[Ineliminable seriousness]\label{l:serious}
Any agent capable of intentional action takes some ends seriously.
\end{lemma}

\begin{proof}
To act intentionally is to pursue an end, which requires devoting some effort and attention to it and treating considerations as bearing on how to achieve it. An agent who took no end seriously in the sense of Definition~\ref{d:serious} would not act, except by reflex. Even the attempt to stop taking ends seriously is the pursuit of an end and is undertaken in the same way. As \citet{Nagel1971} noted, the sceptic's withdrawal from concern is itself conducted with concern.
\end{proof}

\begin{lemma}[Ineliminable detachment]\label{l:detach}
For any reflective agent and any justification she can offer for her seriousness, the detached standpoint can raise the question of what justifies that justification, and no answer is forced on that standpoint.
\end{lemma}

\begin{proof}
A justification of concern for an end cites a consideration that the agent treats as mattering: the end serves another end, or has a value, or is required. From the detached standpoint the agent can ask, of that consideration, why it should matter to her, without presupposing that it does. If she answers by citing a further consideration, the question recurs. If she answers that the consideration has objective worth, the detached standpoint can register the claim but cannot be compelled to endorse it by anything the agent can exhibit, since what she would exhibit is her own conviction, which the step back has placed in question. Nothing in reflection forbids the step back, and a reflective agent is one who can take it.
\end{proof}

The proof of Lemma~\ref{l:detach} does not assume that there is no objective worth. EXI-01 argues that life is meaningful to the extent that one engages with projects of objective worth, and we accept that conclusion. The lemma claims only that the detached standpoint cannot be forced to recognize that worth, so that the regress of justification has no terminus the standpoint must accept. Whether a project is worthwhile is one question; whether its worth can be certified to a standpoint that suspends all of one's evaluative convictions is another, and the absurd concerns the second.

\begin{theorem}\label{t:absurd}
The life of any reflective agent capable of intentional action is absurd in the sense of Definition~\ref{d:absurd}.
\end{theorem}

\begin{proof}
By Lemma~\ref{l:serious} the agent takes some ends seriously. By Lemma~\ref{l:detach}, from the detached standpoint no justification of that seriousness is available that is not open to the same question. By Definition~\ref{d:absurd}, her life is absurd.
\end{proof}

\begin{proposition}[Despair does not fit]\label{p:despair}
Despair is not a fitting response to the absurdity established by Theorem~\ref{t:absurd}.
\end{proposition}

\begin{proof}
By Definition~\ref{d:resp}, despair presents one's ends as worthless or beyond the reach of good. Theorem~\ref{t:absurd} establishes only that no justification of one's seriousness is forced on the detached standpoint. That is the absence of a vindication, not a demonstration of worthlessness: the detached standpoint, which suspends evaluative conviction, is in no better position to certify that one's ends lack worth than to certify that they have it. The situation therefore lacks the feature under which despair presents it, and by Principle~\ref{p:fit} despair does not fit.
\end{proof}

\begin{proposition}[Revolt does not fit]\label{p:revolt}
Revolt is not a fitting response to the absurdity established by Theorem~\ref{t:absurd}.
\end{proposition}

\begin{proof}
Revolt presents the world as an adversary that denies us meaning. The absurd of Definition~\ref{d:absurd} is not a relation between the agent and the world but between two standpoints of the agent. The world plays no part in it, and EXI-06 argues that the universe is indifferent to us, which is not hostility. An indifferent world denies us nothing and cannot be defeated by our persistence. The situation lacks the feature under which revolt presents it, and by Principle~\ref{p:fit} revolt does not fit.
\end{proof}

\begin{theorem}[Irony fits]\label{t:irony}
Irony is a fitting response to the absurdity of life, and it is the only one of the three responses of Definition~\ref{d:resp} that is.
\end{theorem}

\begin{proof}
Irony presents a discrepancy between the seriousness with which one holds one's ends and the support for it available from the most detached standpoint. By Theorem~\ref{t:absurd}, that discrepancy exists: one holds one's ends seriously, and the detached standpoint supplies no support it must accept. Irony also involves continuing to hold one's ends, which Lemma~\ref{l:serious} shows a reflective agent will do in any case and which Proposition~\ref{p:despair} shows the absurd gives no reason to abandon. The situation has the features under which irony presents it, so by Principle~\ref{p:fit} irony fits. By Propositions~\ref{p:despair} and~\ref{p:revolt}, the other two responses do not.
\end{proof}

\begin{remark}\label{r:camus}
Camus's Sisyphus is happy because he is lucid about his condition and refuses to be crushed by it. We grant the lucidity and reject the framing. The punishment of Sisyphus is imposed by gods who intend it as a punishment; his scorn is directed at them and is fitting. The human situation, on Nagel's analysis, contains no such authority. Transposed to it, scorn has no object, and the heroism it confers is borrowed from a myth whose structure our situation lacks.
\end{remark}

\section{Objections and replies}

\begin{objection}[Evasion]\label{o:evasion}
Irony is a way of not taking the absurd seriously. Camus demanded that we keep the contradiction alive and refuse every consolation; irony is a consolation, a smile that lets us return to business as usual.
\end{objection}

\begin{reply}
Irony does not deny the contradiction; it consists in seeing it. What it declines to do is to inflate it into a grievance or a catastrophe. The demand that the absurd be met with constant tension presupposes that the detached standpoint has the authority to make our concerns into a problem that requires heroic management. Lemma~\ref{l:detach} gives it no such authority: it can question but not condemn. To take the absurd exactly as seriously as it deserves is not evasion.
\end{reply}

\begin{objection}[Objective worth dissolves the absurd]\label{o:worth}
If some projects are worthwhile independently of our attitudes, as EXI-01 holds, then the seriousness with which we pursue them is justified, and the life of someone engaged in them is not absurd.
\end{objection}

\begin{reply}
The seriousness is justified, if the projects are worthwhile; the absurd survives because that justification is not one the detached standpoint must accept. Definition~\ref{d:absurd} concerns the availability of a forced justification, not the existence of a sound one. The engaged and worthwhile life is less exposed to despair than the life devoted to trivial ends, since despair is even further from fitting it, but it is no less absurd, and irony fits it as well. This is the sense in which, as the remark at the end of EXI-01-a puts it, the assumption of objective worth denies the premise of the sceptical challenge without refuting it.
\end{reply}

\begin{objection}[Irony corrodes engagement]\label{o:corrode}
Meaning requires active engagement with something worthwhile (EXI-17). The ironist holds her commitments at a distance and is therefore less engaged. Irony is the enemy of meaning.
\end{objection}

\begin{reply}
The objection confuses irony toward one's ends with detachment from them. The irony of Definition~\ref{d:resp} is a second-order attitude toward one's seriousness, and it includes continuing to hold one's ends. A scientist can be absorbed in her research and also aware that, from the standpoint of the universe, nothing requires that anyone find out how proteins fold. The awareness does not diminish the absorption, since the absorption never rested on the universe's endorsement. Ironic noncommitment, the posture of someone who will not care about anything, would corrode engagement; it is a different attitude and one we do not recommend.
\end{reply}

\begin{objection}[Hume's remedy]\label{o:hume}
Hume found that philosophical melancholy dispelled itself when he dined and played backgammon with friends \citep[1.4.7]{Hume1739}. The appropriate response to the absurd is not irony but forgetting.
\end{objection}

\begin{reply}
Forgetting is not a response to the absurd but an interruption of the reflection that discloses it. It is what normally happens, and it is no fault. But the question was what attitude fits the absurd when one does attend to it, and to that question forgetting gives no answer. Irony is the attitude Hume's backgammon player can take when he recalls, between games, what his study showed him.
\end{reply}

\section{What the result does not establish}

The argument does not show that any particular life is meaningful or worthwhile; that question is answered by EXI-01 and EXI-17, and absurdity, as defined here, is shared by meaningful and meaningless lives alike. It does not show that despair is never fitting: a person whose ends are in fact worthless, or who has lost every good, may have reason to despair, but the reason is not the absurd. Nor does it show that irony is obligatory. Principle~\ref{p:fit} concerns fit, and an attitude can fit without being required; someone who never reflects in the way that discloses the absurd has no occasion for irony.

\section{Conclusion}

Life is absurd in Nagel's sense: a reflective agent cannot avoid taking her ends seriously, and cannot obtain from the most detached standpoint a justification of that seriousness that the standpoint must accept. The absurd so understood is the absence of a vindication, not the presence of a refutation. Despair misrepresents it as a verdict of worthlessness, and Camus's revolt misrepresents it as a conflict with a hostile world. Irony presents it accurately, as a discrepancy between commitment and external support, and leaves the commitment in place. It is therefore the appropriate response.

\begin{thebibliography}{9}\small
\bibitem[Camus(1942)]{Camus1942} Camus, A. (1942). \emph{Le Mythe de Sisyphe}. Gallimard. English translation by J. O'Brien, \emph{The Myth of Sisyphus}, Hamish Hamilton, 1955.
\bibitem[Hume(1739)]{Hume1739} Hume, D. (1739). \emph{A Treatise of Human Nature}. Cited by book, part and section.
\bibitem[Metz(2013)]{Metz2013} Metz, T. (2013). \emph{Meaning in Life: An Analytic Study}. Oxford University Press.
\bibitem[Nagel(1971)]{Nagel1971} Nagel, T. (1971). The absurd. \emph{The Journal of Philosophy}, 68(20), 716--727.
\bibitem[Nagel(1986)]{Nagel1986} Nagel, T. (1986). \emph{The View from Nowhere}. Oxford University Press.
\bibitem[Wolf(2010)]{Wolf2010} Wolf, S. (2010). \emph{Meaning in Life and Why It Matters}. Princeton University Press.
\end{thebibliography}

\end{document}
